\documentclass[11pt,a4paper]{article}

\usepackage[utf8]{inputenc}
\usepackage[T1]{fontenc}
\usepackage{amsmath,amssymb}
\usepackage{booktabs,tabularx}
\usepackage{enumitem}
\usepackage{geometry}
\usepackage{hyperref}
\usepackage{xcolor}
\usepackage{listings}
\usepackage{fancyhdr}
\usepackage{titlesec}

\geometry{margin=2.2cm,top=2.8cm,bottom=2.8cm}

\definecolor{codebg}{HTML}{F5F5F5}

\hypersetup{
  colorlinks=true,
  linkcolor=blue!60!black,
  urlcolor=blue!60!black,
  pdftitle={ISLS D3 --- Chat to App},
  pdfauthor={Sebastian Klemm},
}

\lstset{
  basicstyle=\ttfamily\scriptsize,
  breaklines=true,
  frame=single,
  backgroundcolor=\color{codebg},
  keywordstyle=\color{blue!70!black}\bfseries,
  commentstyle=\color{green!50!black}\itshape,
  numbers=left,
  numberstyle=\tiny\color{gray},
  numbersep=5pt,
  tabsize=4,
  showstringspaces=false,
  xleftmargin=1.5em,
  framexleftmargin=1em,
}

\titleformat{\section}{\Large\bfseries}{\thesection}{1em}{}
\titleformat{\subsection}{\large\bfseries}{\thesubsection}{1em}{}
\titleformat{\subsubsection}{\normalsize\bfseries}{\thesubsubsection}{1em}{}

\pagestyle{fancy}
\fancyhf{}
\fancyhead[L]{\small ISLS D3 --- Chat to App}
\fancyhead[R]{\small Klemm, 2026}
\fancyfoot[C]{\thepage}

\begin{document}

% ============================================================
\begin{titlepage}
\begin{center}
\vspace*{1.5cm}
{\Huge\bfseries ISLS D3}\\[0.4cm]
{\LARGE\bfseries Chat to App}\\[0.3cm]
{\LARGE\bfseries Autonome Emission}\\[1.2cm]
{\Large Nat\"urliche Sprache $\to$ TOML $\to$ HDAG $\to$ App}\\[0.3cm]
{\Large Null manuelle Schritte.}\\[1.5cm]
{\large Sebastian Klemm}\\[0.3cm]
{\small 30. M\"arz 2026}\\[2cm]

\begin{tabular}{ll}
\textbf{Repository:} & \texttt{graphity/isls} \\
\textbf{Baut auf:} & D2 (Multi-Domain Proof) \\
\textbf{Crates:} & \texttt{isls-chat} (upgrade), \texttt{isls-cli} (neuer Befehl) \\
\textbf{LLM-Nutzung:} & Intent-Extraktion (1 Call) + Query-Generierung (7-15 Calls) \\
\end{tabular}

\vfill
{\small
D1: TOML $\to$ App. Bewiesen.\\
D2: Drei TOMLs $\to$ drei Apps. Bewiesen.\\
D3: \emph{Ein Satz} $\to$ App. Kein TOML schreiben. Kein Feld definieren.\\
Der User sagt was er will. Das System baut es.\\[0.2cm]
\textbf{Architektur-Prinzip:} Der Chat erzeugt eine TOML. Die TOML
durchl\"auft die bewährte D2-Pipeline. Die deterministische Kette
bleibt intakt. Nur der Eingang wird unscharf.
}
\end{center}
\end{titlepage}

\tableofcontents
\newpage

% ============================================================
\section{Architektur}
% ============================================================

\subsection{Die Kette}

\begin{lstlisting}[language={}]
User: "Ich brauche eine App fuer mein Restaurant
       mit Tischreservierungen und Speisekarte"
    |
    v
[1] LLM-Call: Extrahiere Entities, Fields, Beziehungen
    -> JSON mit Entities + Fields + FKs
    |
    v
[2] JSON -> TOML (deterministisch, kein LLM)
    -> examples/generated.toml
    |
    v
[3] Bestehende D2-Pipeline: forge-v2 --requirements generated.toml
    -> kompilierbare App
    |
    v
[4] Optional: docker-compose up
    -> laufende App im Browser
\end{lstlisting}

\textbf{Kritisch:} Schritt 1 ist der einzige unscharfe Schritt.
Schritte 2--4 sind deterministisch und bewiesen (D1/D2).
Wenn Schritt 1 schlechte Entities extrahiert, wird die App schlecht
aber sie \emph{kompiliert}. Die Pipeline bricht nicht.

\subsection{Warum TOML als Zwischenformat}

Das LLM k\"onnte direkt eine \texttt{AppSpec} erzeugen. Aber:

\begin{enumerate}[nosep]
\item Die TOML ist inspizierbar. Der User kann sie sehen und korrigieren.
\item Die TOML ist reproduzierbar. Derselbe Input erzeugt dieselbe App.
\item Die TOML ist die Schnittstelle die D1/D2 getestet haben.
\item Wenn das LLM Unsinn extrahiert, ist die TOML der Ort wo man
      es sieht --- nicht in 55 generierten Dateien.
\end{enumerate}

% ============================================================
\section{Schritt 1: Entity-Extraktion via LLM}
% ============================================================

\subsection{Der Prompt}

Ein einzelner LLM-Call extrahiert die komplette Anwendungsstruktur:

\begin{lstlisting}[language={}]
You are an expert software architect. Extract the data model
from this application description. Return ONLY valid JSON,
no markdown fences, no explanation.

The JSON must follow this exact schema:
{
  "app_name": "kebab-case-name",
  "description": "one line description",
  "entities": [
    {
      "name": "PascalCaseName",
      "fields": [
        {
          "name": "snake_case_name",
          "type": "String|i32|i64|f64|bool",
          "nullable": false,
          "unique": false,
          "default": null
        }
      ],
      "foreign_keys": [
        {
          "target": "OtherEntityName",
          "nullable": false
        }
      ]
    }
  ]
}

Rules:
- Every app needs a User entity with: email (String, unique),
  password_hash (String), role (String, default "user"),
  is_active (bool, default true).
- Use i64 for money (cents), i32 for quantities, String for text,
  bool for flags, f64 for measurements.
- Add created_at/updated_at automatically - do NOT include them.
- id is automatic - do NOT include it.
- Foreign keys reference other entities by PascalCase name.
- If an entity "belongs to" another, add a foreign_key.
- Name entities in PascalCase singular: "Product", not "products".
- Name fields in snake_case: "unit_price", not "unitPrice".
- Include 4-12 entities. Not less, not more.
- Every entity needs at least 1 own field besides foreign keys.

User's description:
{user_message}
\end{lstlisting}

\subsection{JSON zu TOML Konvertierung}

Deterministisch. Kein LLM.

\begin{lstlisting}[language=C]
/// Convert LLM-extracted JSON to a valid TOML string.
pub fn json_to_toml(json: &serde_json::Value) -> Result<String> {
    let mut toml = String::new();

    // [app] section
    toml.push_str("[app]\n");
    toml.push_str(&format!("name = \"{}\"\n",
        json["app_name"].as_str().unwrap_or("my-app")));
    toml.push_str(&format!("description = \"{}\"\n\n",
        json["description"].as_str().unwrap_or("Generated application")));

    // [backend] section
    toml.push_str("[backend]\n");
    toml.push_str("language = \"rust\"\n");
    toml.push_str("framework = \"actix-web\"\n");
    toml.push_str("database = \"postgresql\"\n");
    toml.push_str("auth_method = \"jwt\"\n\n");

    // [frontend] section
    toml.push_str("[frontend]\n");
    toml.push_str("type = \"spa\"\n");
    toml.push_str("framework = \"vanilla\"\n");
    toml.push_str("styling = \"minimal\"\n\n");

    // [deployment] section
    toml.push_str("[deployment]\n");
    toml.push_str("containerized = true\n");
    toml.push_str("compose = true\n\n");

    // [[entities]] sections
    if let Some(entities) = json["entities"].as_array() {
        for entity in entities {
            toml.push_str("[[entities]]\n");
            toml.push_str(&format!("name = \"{}\"\n",
                entity["name"].as_str().unwrap_or("Unknown")));

            // Fields
            if let Some(fields) = entity["fields"].as_array() {
                toml.push_str("fields = [\n");
                for field in fields {
                    toml.push_str(&format!(
                        "    {{ name = \"{}\", type = \"{}\"",
                        field["name"].as_str().unwrap_or("field"),
                        field["type"].as_str().unwrap_or("String"),
                    ));
                    if field["nullable"].as_bool().unwrap_or(false) {
                        toml.push_str(", nullable = true");
                    }
                    if field["unique"].as_bool().unwrap_or(false) {
                        toml.push_str(", unique = true");
                    }
                    if let Some(def) = field["default"].as_str() {
                        toml.push_str(&format!(", default = \"{}\"", def));
                    }
                    toml.push_str(" },\n");
                }
                toml.push_str("]\n");
            }

            // Foreign keys
            if let Some(fks) = entity["foreign_keys"].as_array() {
                if !fks.is_empty() {
                    toml.push_str("foreign_keys = [\n");
                    for fk in fks {
                        toml.push_str(&format!(
                            "    {{ target = \"{}\"",
                            fk["target"].as_str().unwrap_or("Unknown"),
                        ));
                        if fk["nullable"].as_bool().unwrap_or(false) {
                            toml.push_str(", nullable = true");
                        }
                        toml.push_str(" },\n");
                    }
                    toml.push_str("]\n");
                }
            }
            toml.push_str("\n");
        }
    }

    Ok(toml)
}
\end{lstlisting}

\subsection{Validierung des LLM-Outputs}

Bevor die TOML an die Pipeline geht:

\begin{lstlisting}[language=C]
pub fn validate_extracted_spec(json: &serde_json::Value) -> Result<()> {
    let entities = json["entities"].as_array()
        .ok_or_else(|| anyhow!("no entities array"))?;

    // Must have entities
    if entities.is_empty() {
        return Err(anyhow!("no entities extracted"));
    }

    // Must have a User entity
    let has_user = entities.iter().any(|e|
        e["name"].as_str() == Some("User"));
    if !has_user {
        return Err(anyhow!("no User entity --- required for auth"));
    }

    // Each entity must have a name
    for entity in entities {
        let name = entity["name"].as_str()
            .ok_or_else(|| anyhow!("entity without name"))?;

        // PascalCase check
        if name.chars().next().map_or(true, |c| !c.is_uppercase()) {
            return Err(anyhow!("entity '{}' not PascalCase", name));
        }

        // Must have fields or foreign_keys (or both)
        let field_count = entity["fields"].as_array()
            .map_or(0, |f| f.len());
        let fk_count = entity["foreign_keys"].as_array()
            .map_or(0, |f| f.len());
        if field_count == 0 && fk_count == 0 {
            return Err(anyhow!("entity '{}' has no fields and no FKs", name));
        }

        // FK targets must reference existing entities
        if let Some(fks) = entity["foreign_keys"].as_array() {
            let entity_names: Vec<&str> = entities.iter()
                .filter_map(|e| e["name"].as_str())
                .collect();
            for fk in fks {
                let target = fk["target"].as_str()
                    .ok_or_else(|| anyhow!("FK without target"))?;
                if !entity_names.contains(&target) {
                    return Err(anyhow!(
                        "FK target '{}' not in entities", target));
                }
            }
        }

        // Field types must be valid
        if let Some(fields) = entity["fields"].as_array() {
            for field in fields {
                let ft = field["type"].as_str().unwrap_or("");
                if !["String", "i32", "i64", "f64", "bool"].contains(&ft) {
                    return Err(anyhow!(
                        "invalid field type '{}' in {}", ft, name));
                }
            }
        }
    }

    Ok(())
}
\end{lstlisting}

Wenn die Validierung fehlschl\"agt: dem User die Fehlermeldung zeigen
und erneut fragen, nicht die Pipeline starten.

% ============================================================
\section{Neuer CLI-Befehl: \texttt{forge-chat}}
% ============================================================

\begin{lstlisting}[language=bash]
# Chat-to-App: ein Satz rein, App raus
isls forge-chat \
    --message "Restaurant app with table reservations and menu" \
    --api-key $OPENAI_API_KEY \
    --output ./output/restaurant

# Was passiert:
# 1. LLM extrahiert Entities -> JSON
# 2. JSON -> TOML (saved as output/restaurant.toml)
# 3. forge-v2 --requirements output/restaurant.toml
# 4. Output in ./output/restaurant/
\end{lstlisting}

\subsection{Implementation in \texttt{isls-cli}}

\begin{lstlisting}[language=C]
/// The forge-chat subcommand
pub fn forge_chat(
    message: &str,
    api_key: &str,
    model: &str,        // default: "gpt-4o"
    output_dir: &Path,
) -> Result<()> {
    // 1. Extract entities via LLM
    tracing::info!("[Chat] Extracting entities from description...");
    let oracle = OpenAiOracle::new(api_key.into())
        .with_model(model);
    let prompt = build_extraction_prompt(message);
    let json_str = oracle.call(&prompt, 4096)?;

    // 2. Parse and validate JSON
    let json: serde_json::Value = serde_json::from_str(&json_str)
        .map_err(|e| anyhow!("LLM returned invalid JSON: {}", e))?;
    validate_extracted_spec(&json)?;
    tracing::info!("[Chat] Extracted {} entities",
        json["entities"].as_array().map_or(0, |e| e.len()));

    // 3. Convert to TOML
    let toml_content = json_to_toml(&json)?;
    let toml_path = output_dir.join("spec.toml");
    std::fs::create_dir_all(output_dir)?;
    std::fs::write(&toml_path, &toml_content)?;
    tracing::info!("[Chat] TOML written to {}", toml_path.display());

    // 4. Print extracted spec for user review
    println!("\n--- Extracted Specification ---");
    println!("{}", toml_content);
    println!("--- End Specification ---\n");

    // 5. Run the proven D2 pipeline
    tracing::info!("[Chat] Starting forge pipeline...");
    forge_v2(&toml_path, Some(api_key), model, output_dir)?;

    tracing::info!("[Chat] Complete. Output in {}", output_dir.display());
    Ok(())
}
\end{lstlisting}

\subsection{CLI Argument Parsing}

\begin{lstlisting}[language=C]
// In main.rs, add to the match:
"forge-chat" => {
    let message = args.get("--message")
        .or_else(|| args.get("-m"))
        .ok_or_else(|| anyhow!("--message required"))?;
    let api_key = args.get("--api-key")
        .or_else(|| std::env::var("OPENAI_API_KEY").ok().as_ref())
        .ok_or_else(|| anyhow!("--api-key or OPENAI_API_KEY required"))?;
    let model = args.get("--model").unwrap_or(&"gpt-4o".into());
    let output = args.get("--output")
        .ok_or_else(|| anyhow!("--output required"))?;
    forge_chat(message, api_key, model, Path::new(output))?;
}
\end{lstlisting}

% ============================================================
\section{Wo der Code hinkommt}
% ============================================================

\begin{center}
\small
\begin{tabular}{l l p{7cm}}
\toprule
\textbf{Funktion} & \textbf{Datei} & \textbf{Beschreibung} \\
\midrule
\texttt{build\_extraction\_prompt} & \texttt{isls-chat/src/lib.rs} & Baut den LLM-Prompt f\"ur Entity-Extraktion \\
\texttt{validate\_extracted\_spec} & \texttt{isls-chat/src/lib.rs} & Validiert den JSON-Output \\
\texttt{json\_to\_toml} & \texttt{isls-chat/src/lib.rs} & Konvertiert JSON zu TOML \\
\texttt{forge\_chat} & \texttt{isls-cli/src/main.rs} & Orchestriert den gesamten Flow \\
\bottomrule
\end{tabular}
\end{center}

\texttt{isls-chat} existiert bereits als Crate. Es hat Intent-Erkennung
via Keywords. Die LLM-basierte Entity-Extraktion und JSON-zu-TOML
Konvertierung werden hinzugef\"ugt.

\texttt{isls-cli} bekommt den neuen \texttt{forge-chat} Befehl der
\texttt{isls-chat} f\"ur Extraktion und dann \texttt{forge-v2} f\"ur
Generierung aufruft.

% ============================================================
\section{Testf\"alle}
% ============================================================

\subsection{Test 1: Restaurant}

\begin{lstlisting}[language=bash]
isls forge-chat \
    --message "Restaurant app with table reservations,
               menu with categories, and customer reviews" \
    --api-key %OPENAI_API_KEY% \
    --output ./output/restaurant
cd output/restaurant/backend && cargo build && cd ../../..
\end{lstlisting}

Erwartete Entities: User, Restaurant (oder Menu), Category, MenuItem,
Table, Reservation, Review. 6--8 Entities.

\subsection{Test 2: Bibliothek}

\begin{lstlisting}[language=bash]
isls forge-chat \
    --message "Library management: books, members,
               borrowing and returns, overdue tracking" \
    --api-key %OPENAI_API_KEY% \
    --output ./output/library
cd output/library/backend && cargo build && cd ../../..
\end{lstlisting}

Erwartete Entities: User, Book, Author, Member, Loan, Category. 5--7 Entities.

\subsection{Test 3: Fitness-Studio}

\begin{lstlisting}[language=bash]
isls forge-chat \
    --message "Gym management with members, trainers,
               class schedules, and membership plans" \
    --api-key %OPENAI_API_KEY% \
    --output ./output/gym
cd output/gym/backend && cargo build && cd ../../..
\end{lstlisting}

Erwartete Entities: User, Member, Trainer, Class, Schedule,
MembershipPlan, Booking. 6--8 Entities.

\subsection{Akzeptanzkriterium}

Alle drei Testf\"alle m\"ussen:
\begin{enumerate}[nosep]
\item Entities extrahieren (JSON valide, Validierung bestanden)
\item TOML erzeugen (inspizierbar, korrekte Syntax)
\item \texttt{cargo build} bestehen (0 Errors)
\item Bonus: Docker starten + Login funktionieren
\end{enumerate}

Es ist \emph{akzeptabel} wenn die extrahierten Entities nicht perfekt
sind --- das LLM interpretiert ``Restaurant'' vielleicht anders als
erwartet. Solange es kompiliert und CRUD funktioniert, ist D3 bestanden.
Die Qualit\"at der Extraktion verbessert sich mit besseren Prompts
und sp\"ater mit Norm-Matching (D4).

% ============================================================
\section{Was D3 NICHT ist}
% ============================================================

\begin{itemize}[nosep]
\item \textbf{Kein Chat-Dialog.} D3 ist ein Single-Shot: ein Satz rein,
      eine App raus. Iteratives Verfeinern (``f\"uge ein Gewichtsfeld
      hinzu'') ist D4+.
\item \textbf{Kein Studio.} D3 ist CLI-only. Die Web-UI kommt sp\"ater.
\item \textbf{Kein Norm-Matching.} D3 benutzt das LLM direkt f\"ur
      Entity-Extraktion. Norm-basierte Extraktion ist D4.
\item \textbf{Kein Frontend-Customizing.} Das Frontend ist dasselbe
      Skelett wie bei D1/D2. Besser wird es wenn die Frontend-Generatoren
      verbessert werden.
\end{itemize}

% ============================================================
\section{Verifikation}
% ============================================================

\begin{lstlisting}[language=bash]
# 1. Workspace
cargo build --workspace --release
cargo test --workspace

# 2. D1/D2 Regression
rm -rf output/warehouse
cargo run --release -p isls-cli -- forge-v2 \
    --requirements examples/warehouse.toml \
    --mock-oracle --output ./output/warehouse
cd output/warehouse/backend && cargo build && cd ../../..

# 3. D3: Restaurant
cargo run --release -p isls-cli -- forge-chat \
    --message "Restaurant app with table reservations, menu management, and customer reviews" \
    --api-key %OPENAI_API_KEY% \
    --output ./output/restaurant
cd output/restaurant/backend && cargo build && cd ../../..

# 4. D3: Library
cargo run --release -p isls-cli -- forge-chat \
    --message "Library system for books, members, borrowing and returns" \
    --api-key %OPENAI_API_KEY% \
    --output ./output/library
cd output/library/backend && cargo build && cd ../../..

# 5. D3: Gym
cargo run --release -p isls-cli -- forge-chat \
    --message "Gym management with members, trainers, classes and memberships" \
    --api-key %OPENAI_API_KEY% \
    --output ./output/gym
cd output/gym/backend && cargo build && cd ../../..

# Alle muessen cargo build bestehen.
\end{lstlisting}

% ============================================================
\section{Constraints f\"ur Claude Code}
% ============================================================

\begin{enumerate}
\item \textbf{Der Extraction-Prompt ist exakt wie in \S2.1 definiert.}
      Nicht vereinfachen. Die Qualit\"at der Extraktion h\"angt vom
      Prompt ab.
\item \textbf{JSON $\to$ TOML ist deterministisch.} Kein LLM-Call.
      Reine String-Formatierung.
\item \textbf{Die erzeugte TOML wird gespeichert} als
      \texttt{output/\{name\}/spec.toml} f\"ur sp\"atere Inspektion.
\item \textbf{Die erzeugte TOML wird auf stdout gedruckt} damit
      der User sie sehen kann bevor die Pipeline startet.
\item \textbf{Validierung vor Pipeline.} Wenn das LLM ung\"ultiges
      JSON oder fehlende Entities liefert: Fehler ausgeben, nicht
      die Pipeline starten.
\item \textbf{D1/D2 darf nicht brechen.} \texttt{forge-v2} bleibt
      unver\"andert. \texttt{forge-chat} ruft \texttt{forge-v2} auf.
\item \textbf{Ein LLM-Call f\"ur Extraktion.} Nicht mehrere Runden.
      Ein Prompt, ein Response, eine TOML.
\item \textbf{Default-Modell: gpt-4o.}
\item \textbf{Alle bestehenden Tests m\"ussen bestehen.}
\item \textbf{Neuer Test:} Unit-Test f\"ur \texttt{json\_to\_toml} ---
      bekanntes JSON rein, erwartete TOML raus.
\item \textbf{Neuer Test:} Unit-Test f\"ur \texttt{validate\_extracted\_spec}
      --- g\"ultiges JSON akzeptiert, ung\"ultiges abgelehnt.
\end{enumerate}

% ============================================================
\section{Was D3 beweist}
% ============================================================

Wenn alle drei Testf\"alle kompilieren:

\begin{itemize}[nosep]
\item Ein Satz nat\"urliche Sprache reicht um eine funktionief\"ahige
      Anwendung zu erzeugen.
\item Die D2-Pipeline ist robust genug um LLM-extrahierte Specs
      zu verarbeiten --- nicht nur handgeschriebene TOMLs.
\item Die Kette Chat $\to$ JSON $\to$ TOML $\to$ HDAG $\to$ App
      funktioniert end-to-end.
\item Der Generator ist nicht nur dom\"anenunabh\"angig (D2) sondern
      auch inputunabh\"angig: TOML oder nat\"urliche Sprache,
      dasselbe Ergebnis.
\end{itemize}

D3 ist der Moment ab dem Homer den Generator benutzen kann.
Er muss keine TOML schreiben. Er muss kein Entity-Modell entwerfen.
Er sagt was er will. Der Rest passiert.

\vfill
\noindent\rule{\textwidth}{0.4pt}
\begin{center}
\small
ISLS D3 --- Chat to App --- Sebastian Klemm --- 30. M\"arz 2026\\
Ein Satz. Eine App. Null manuelle Schritte.\\
Die Trompete h\"ort jetzt auf die Stimme.\\
\url{https://github.com/LashSesh/graphity}
\end{center}

\end{document}
